\documentclass[12pt]{article}
\usepackage{setspace}
\onehalfspacing

\usepackage{iftex}
\ifPDFTeX
  \usepackage[utf8]{inputenc}
  \usepackage[T1]{fontenc}
\else
  \usepackage{fontspec}
\fi
\usepackage[spanish, es-tabla]{babel}
\usepackage[margin=1in]{geometry}
\usepackage{enumitem}
\usepackage{hyperref}
\usepackage{microtype}
\usepackage{biblatex}
\addbibresource{refs.bib}

\hypersetup{
  colorlinks=true,
  linkcolor=black,
  urlcolor=blue,
  citecolor=black
}

\setlist[itemize]{leftmargin=*, itemsep=0.35em}
\newcommand{\institucion}{Universidad de Costa Rica}
\newcommand{\campus}{Ciudad Universitaria Rodrigo Facio}
\newcommand{\facultad}{Facultad de Ciencias Básicas}
\newcommand{\escuela}{Escuela de Física}
\newcommand{\curso}{Mecánica Estadística Computacional}
\newcommand{\tituloanteproyecto}{Proyecto de Investigación}
\newcommand{\titulo}{Evaluación de Predicciones Monte Carlo para Analizar la Efectividad de Medidas contra la Propagación de COVID-19}
\newcommand{\autor}{Gustavo A. Calvo Gutiérrez}
\newcommand{\lugarfecha}{San José, 2026.}

\begin{document}

\begin{titlepage}
  \begin{center}
    \LARGE \textbf{\textsc{\institucion}}
    \Large
    \\[1cm]
    \textsc{\campus}
    \\[0.1cm]
    \textsc{\facultad}
    \\[0.1cm]
    \textsc{\escuela} \\
    \vspace*{20 mm}
    \LARGE{\textsc{\textbf{\curso}}}
  \end{center}

  \vspace{2cm}

  \hrulefill
  \begin{center}
    \large \textbf{\textsc{\tituloanteproyecto}}
    \\[0.5cm]
    \textsc{\titulo}
  \end{center}
  \hrulefill

  \vfill

  \begin{center}
    \large \textbf{\autor}
  \end{center}

  \vspace{1cm}

  \begin{center}
    \textit{\lugarfecha}
  \end{center}
\end{titlepage}

\section{Pregunta de Investigación y Justificación}

\subsection{Pregunta de Investigación}

\noindent
\textit{
  ¿En qué medida las predicciones obtenidas mediante un modelo de simulación
  Monte Carlo basado en el procedimiento propuesto por Xie \cite{Xie2020}, al
  contrastarse con datos reales de contagio, recuperación, densidad poblacional
  efectiva y cronología de medidas por zona y por periodo, permiten identificar
  qué medidas se asocian con una mayor reducción de la transmisión de COVID-19?
}

\subsection{Justificación}

El estudio de la propagación de COVID-19 continúa siendo relevante como caso de
referencia para comprender fenómenos epidémicos donde la transmisión depende de
interacciones sociales, cambios de comportamiento y políticas públicas aplicadas
en distintos momentos. Los modelos estocásticos ofrecen una ventaja importante
frente a aproximaciones puramente deterministas, ya que permiten representar
explícitamente la incertidumbre asociada a la ocurrencia de nuevos contagios y a
la variabilidad temporal de la transmisión.

El trabajo de Xie \cite{Xie2020} plantea una estrategia basada en simulación
Monte Carlo para describir la dinámica de contagio mediante tasas de infección
variables en el tiempo, junto con distribuciones probabilísticas para modelar
tanto la cantidad de personas contagiadas por cada caso activo como el momento
en que dichos contagios ocurren.

El presente anteproyecto propone tomar ese artículo como punto de partida
metodológico para comparar predicciones simuladas con datos reales observados en
zonas específicas. La meta no es únicamente reproducir un modelo existente,
sino utilizarlo como herramienta analítica para explorar qué medidas se asocian
con reducciones más claras en la velocidad de transmisión. Este enfoque es
pertinente porque permite vincular la simulación con evidencia empírica. Para
ello, el estudio puede apoyarse en fuentes comparables de series epidemiológicas
y de registro de políticas públicas, como \textit{Our World in Data}
\cite{owid-coronavirus} y el
\textit{Oxford COVID-19 Government Response Tracker} \cite{Hale2021}.

Al poner en relación las curvas generadas por el modelo con la evolución real de
los casos y con la cronología de las medidas aplicadas, se podrá analizar qué
intervenciones o combinaciones de intervenciones se asocian con cambios
relevantes en la transmisión.

\newpage

\section{Objetivos}

\subsection{Objetivo General}

Evaluar la capacidad de un modelo de simulación Monte Carlo, basado
en el procedimiento propuesto por Xie \cite{Xie2020}, para comparar
predicciones epidemiológicas con datos reales de COVID-19 y analizar la
asociación entre distintas medidas y la reducción de la transmisión en zonas
específicas.

\subsection{Objetivos Específicos}

\begin{itemize}
  \item Adaptar el procedimiento Monte Carlo de Xie \cite{Xie2020} a un
        conjunto de datos temporales de COVID-19 organizado por zona geográfica.
  \item Recopilar y estructurar series de tiempo con variables relevantes,
        tales como personas infectadas, recuperadas, casos activos, densidad
        poblacional efectiva y momentos de aplicación de medidas documentadas
        en fuentes comparables \cite{owid-coronavirus,Hale2021}.
  \item Modelar la probabilidad diaria de contagio como un componente aleatorio
        dependiente del tiempo y sensible a cambios en las condiciones de
        movilidad, contacto social e implementación de medidas.
  \item Comparar las trayectorias simuladas con los datos reales observados
        para identificar discrepancias, patrones recurrentes y posibles ajustes
        del modelo.
  \item Analizar, mediante escenarios comparativos, qué medidas o
        combinaciones de medidas parecen asociarse con una mayor reducción de
        la transmisión a lo largo del tiempo.
\end{itemize}

\newpage

\section{Marco Teórico}

La simulación de epidemias puede abordarse desde modelos deterministas, como
los de tipo SIR, o desde modelos estocásticos que incorporan explícitamente la
incertidumbre del proceso de contagio.
Los modelos estocásticos consideran que la transmisión
ocurre a través de eventos aleatorios cuya realización concreta puede variar
incluso bajo condiciones iniciales similares. Esta segunda perspectiva resulta
especialmente útil cuando se desea representar fluctuaciones diarias,
heterogeneidad espacial y cambios abruptos en el comportamiento social.

En el trabajo de Xie \cite{Xie2020},
la propuesta modela la propagación de COVID-19 como un proceso de puntos en el
que cada persona infecciosa puede generar un número aleatorio de nuevos
contagios. La cantidad de infecciones secundarias se representa mediante una
distribución de Poisson de una tasa de infección \(R_t\),
mientras que el momento en que ocurre cada nuevo contagio se modela mediante
una distribución binomial negativa con media \(\mu_T\). Esta construcción
permite simular tanto la cantidad como el tiempo de aparición de los nuevos
casos.

Un punto central del modelo es que \(R_t\) no se considera constante. La tasa
de transmisión cambia con el tiempo para reflejar variaciones en el contacto
social, la respuesta institucional y el comportamiento de la población. En el
artículo de referencia, la selección de patrones para \(R_t\) combina ajuste a
datos observados con interpretación de los efectos de intervenciones públicas
\cite{Xie2020}. Esta idea es especialmente valiosa para el presente proyecto,
porque permite relacionar cambios en la intensidad de la transmisión con la
introducción, el retiro o la combinación de distintas medidas.

Desde el punto de vista empírico, el análisis de medidas requiere series de
tiempo comparables. No basta con observar el número acumulado de casos; es
necesario considerar variables que cambian día con día, tales como el número de
personas activamente infectadas, recuperadas, la población potencialmente
expuesta y la densidad poblacional efectiva en contextos específicos. Aunque la
densidad poblacional total de una zona es relativamente estable, la densidad
efectiva de interacción puede variar con restricciones de movilidad, cierres
parciales, campañas de vacunación y cambios en los patrones de circulación.
Para documentar estos factores, resultan particularmente útiles bases de datos
como \textit{Our World in Data} \cite{owid-coronavirus}, que reúne indicadores
epidemiológicos comparables, y el
\textit{Oxford COVID-19 Government Response Tracker} \cite{Hale2021},
que codifica la evolución de distintas políticas públicas.

El proyecto también se apoya en la idea de validación por contraste entre
simulación y observación.
Si el modelo captura adecuadamente cambios en la tendencia de contagios
tras la introducción de ciertas medidas, entonces puede utilizarse para
explorar escenarios alternativos y discutir su efectividad relativa. Si no lo
hace, las discrepancias mismas aportan información sobre variables omitidas o
supuestos que deben refinarse. En consecuencia, la interpretación de resultados
se planteará en términos de asociación y consistencia empírica, no como una
demostración causal concluyente.

\newpage

\section{Aplicación de la Técnica de Monte Carlo}

\subsection{Recopilación y Organización de Datos}

La primera etapa consistirá en construir una base de datos temporal para zonas
específicas de análisis. Se buscará registrar, para cada unidad espacial y cada
fecha, variables como casos nuevos, casos activos, personas recuperadas,
fallecimientos, población total o estimada, densidad poblacional efectiva y
periodos en que entraron en vigor distintas medidas. Entre estas medidas pueden
incluirse cuarentenas, obligatoriedad de mascarillas, restricciones de aforo,
campañas de vacunación u otras intervenciones registradas sistemáticamente en
fuentes como \textit{Our World in Data} \cite{owid-coronavirus} y el
\textit{Oxford COVID-19 Government Response Tracker} \cite{Hale2021}. La
organización en series de tiempo permitirá observar la evolución del problema y
comparar periodos anteriores y posteriores a cada intervención.

\subsection{Construcción del Modelo Estocástico}

Tomando como referencia el trabajo de Xie \cite{Xie2020}, se propondrá un
modelo en el que cada persona infectada tenga una probabilidad de contagiar a
otras personas en cada día del periodo de estudio. La parte aleatoria del
problema se representará precisamente en esa posibilidad diaria de transmisión.
En lugar de asumir que todos los contagios siguen una trayectoria fija, se
permitirá que el número de infecciones nuevas surja de distribuciones
probabilísticas calibradas con los datos observados.

De manera conceptual, el procedimiento será el siguiente:

\begin{itemize}
  \item Se parte de un número inicial de personas infectadas en una zona y
        fecha determinadas.
  \item Para cada día, se estima una tasa de transmisión asociada al contexto
        del periodo, la cual puede disminuir o aumentar según las medidas
        vigentes, la dinámica social observada y otras condiciones relevantes
        documentadas en los datos.
  \item Cada caso activo puede generar nuevos contagios de forma aleatoria, y
        el momento exacto de esos contagios también se modela
        probabilísticamente.
  \item El proceso se repite durante toda la ventana temporal de estudio para
        obtener trayectorias simuladas de casos nuevos, casos activos y
        acumulados.
\end{itemize}

Este enfoque es coherente con Monte Carlo porque el comportamiento del sistema
se explora a través de muchas realizaciones posibles del mismo proceso
aleatorio. Así, en vez de producir una sola predicción puntual, se obtiene una
distribución de resultados plausibles.

\subsection{Comparación con Datos Reales}

Una vez generadas múltiples trayectorias simuladas, estas se compararán con las
series reales de cada zona. La comparación permitirá examinar si el modelo
logra reproducir momentos de crecimiento acelerado, picos de contagio, periodos
de estabilización y descensos posteriores. También permitirá medir si las
modificaciones temporales de la tasa de transmisión, introducidas para
representar políticas sanitarias concretas, son consistentes con los cambios
efectivamente observados en los datos.

La utilidad principal de esta comparación radica en que las medidas pueden
interpretarse como factores que alteran la estructura de contacto entre
personas o la probabilidad de que un contacto resulte en contagio. Si, por
ejemplo, una reducción sostenida en la transmisión coincide con periodos de
mayor uso de mascarillas, restricciones más estrictas de movilidad o aumentos
en la cobertura de vacunación, el modelo puede capturar ese cambio mediante una
disminución de la tasa de infección y contrastarlo con la realidad observada.

\subsection{Evaluación de la Efectividad de Medidas}

Finalmente, el proyecto comparará distintos escenarios temporales y espaciales
para determinar cuáles medidas parecen asociarse con mejores resultados. No se
asumirá de entrada que una sola intervención explica toda la dinámica del
brote; más bien, se analizará cómo distintas medidas, tomadas de forma
individual o combinada, modifican la evolución de la curva epidémica. La
interpretación de resultados se apoyará en la proximidad entre las simulaciones
y los datos reales, así como en la consistencia del comportamiento estimado de
la tasa de transmisión a lo largo del tiempo.

Este procedimiento permitirá construir una evaluación comparativa con
fundamento cuantitativo, donde la simulación Monte Carlo funcione como puente
entre la formulación teórica del contagio y la evidencia empírica disponible.

\newpage

\printbibliography

\end{document}
